\chapter{Symbolic Language of Resonance}

%==============================================================
%  CHAPTER 4 : Symbolic Language of Resonance
%==============================================================

\begin{center}
\emph{“If motion is the universe’s first verb, a single oscillation is its smallest syllable.”}
\end{center}

\bigskip
We begin by fixing the \emph{alphabet}. A single, self-consistent oscillation \(O_i\) is the irreducible letter; all higher constructs are words and sentences made from these letters.

%--------------------------------------------------------------------
\section{Alphabet of Primitives}

\begin{table}[H]
  \centering
  \caption{RCM alphabet of primitive symbols and their roles. Types refer to the four sorts in the typed grammar.}
  \label{tab:rcm-alphabet}
  \renewcommand{\arraystretch}{1.2}
  \begin{tabular}{@{} >{$}l<{$} >{\scshape}p{0.18\linewidth} >{\raggedright\arraybackslash}p{0.62\linewidth} @{}}
    \toprule
    \multicolumn{1}{c}{\textbf{Symbol}} & \textbf{Type} & \textbf{Meaning / Role} \\
    \midrule
    O_i = [A_i,\varsigma_i,\omega_i] & Term &
      \textbf{Primitive oscillation}. Triple of amplitude \(A_i\), phase \(\varsigma_i\), and frequency \(\omega_i\). \\

    \RB{\omega} & Bundle &
      \textbf{Resonance bundle}. Coherent set of oscillations centered at frequency \(\omega\) with density \(\rho_{\omega}(O)\), normalized as \(\int \rho_{\omega}(O)\,dO = 1\). \\

    K_{\omega\varepsilon}(\tau) & Kernel &
      \textbf{Memory kernel}. Delayed (causal) coupling from bundle \(\varepsilon\) to bundle \(\omega\) at lag \(\tau\). \\

    p(\vartheta) & Scalar field &
      \textbf{Pitch field}. \(p = d(\ln r)/d\vartheta\); acts as a Weyl-like gauge field on spiral geometry. \\

    \Composer\{\,\cdot\,\} & Bundle \(\times\) Bundle \(\to\) Bundle &
      \textbf{Composer}. Joins bundles nose-to-tail when boundary pitches match (Axiom C-1) and is associative (Axiom C-2). \\

    T_{\vartheta} & Operator &
      \textbf{Scaling operator}. \((A,\omega)\mapsto (e^{\vartheta}A,\; \omega e^{-\vartheta})\); global pitch rescaling \(\Leftrightarrow\) Lorentz-like boost. \\

    \hat{\#},\ \hat{D},\ \hat{\$}_n,\ \hat{M}_{\tau} & Operators &
      Phase shift, pitch derivative, winding-\(n\) projector, and memory filter, respectively. \\

    (O) & Integer &
      \textbf{Topological charge} of a single oscillation. \\

    \langle \cdot \rangle & Functional &
      Pitch-period average: \(\displaystyle \langle X\rangle=\frac{1}{2\pi}\int_{\vartheta}^{\vartheta+2\pi} X(\vartheta')\,d\vartheta'\). \\
    \bottomrule
  \end{tabular}
\end{table}



%--------------------------------------------------------------------
\section{Typing Rules and Bundle Orders}

\begin{itemize}[leftmargin=1.5em]
  \item \textbf{Terminals} — order-0 objects \(O_i\).
  \item \textbf{Bundle\(\langle k\rangle\)} — a coherent set of Bundle\(\langle k-1\rangle\) objects. Example: an entangled photon pair is Bundle\(\langle 1\rangle\).
  \item \textbf{Kernel} — maps Bundle\(\langle m\rangle\) \(\times\) time \(\to\) Bundle\(\langle n\rangle\).
  \item \textbf{Operator} — endomorphism on a fixed Bundle\(\langle k\rangle\).
  \item \textbf{Statement} — an equation or transformation type-checked by the grammar.
\end{itemize}

\noindent\textit{Normal-form theorem.} Every finite expression reduces uniquely to a
canonical sequence “projector \(\to\) derivative \(\to\) phase-shift \(\to\) explicit kernel convolution.”

%--------------------------------------------------------------------
\section{Composer Axioms}

\textbf{Axiom C-1 (Pitch matching).}\;
\(\Composer\{\RB{1},\RB{2}\}\) exists iff boundary pitches satisfy
\(p_{\mathrm{out}}(\RB{1})=p_{\mathrm{in}}(\RB{2})\).

\medskip
\textbf{Axiom C-2 (Associativity).}\;
If defined,
\(\Composer\{\Composer\{\RB{1},\RB{2}\},\RB{3}\}=\Composer\{\RB{1},\Composer\{\RB{2},\RB{3}\}\}\).

%--------------------------------------------------------------------
\section{Operator Algebra Snapshot}

\[
[\hat{\#},\hat{D}]=0,\qquad
[\hat{D},\hat{M}_{\tau}]\neq 0,\qquad
\hat{\$}_n^{\,2}=\hat{\$}_n.
\]

%--------------------------------------------------------------------
\section{Worked Example — Entangled Photon Pair}

\begin{enumerate}[leftmargin=1.5em]
  \item Two primitives \(O_{1},O_{2}\) at \(800\,\mathrm{nm}\)
        \(\Rightarrow\) bundles \(\RB{1},\RB{2}\).
  \item Apply opposite phase shifts:
        \(\hat{\#}_{\pi/4}\RB{1},\ \hat{\#}_{-\,\pi/4}\RB{2}\).
  \item Compose:
        \(\Composer\{\hat{\#}_{\pi/4}\RB{1},\hat{\#}_{-\,\pi/4}\RB{2}\}\)
        (Bundle\(\langle 1\rangle\)).
  \item Memory-filter in ring cavity:
        \(\hat{M}_{10\,\mathrm{ps}}\).
  \item Project winding-1 component:
        \(\hat{\$}_{1}\,\hat{M}_{10\,\mathrm{ps}}\,\Composer\{\ldots\}\).
\end{enumerate}

%--------------------------------------------------------------------
\section{Why This Matters}

\begin{itemize}[leftmargin=1.5em]
  \item \textbf{Compression} — Convert PDE-heavy models into compact symbolic strings.
  \item \textbf{Composability} — Any two \RCM\ sentences glue via \(\Composer\{\cdot\}\) as long as pitches match.
  \item \textbf{Automated reasoning} — Normal forms enable theorem-prover checks of conservation laws or bandwidth limits.
\end{itemize}

\bigskip
\textbf{Next:} Section~\ref{sec:composer-operator} exploits \(\Composer\{\cdot\}\) to build spiral molecules
and higher-order bundles, while Section~\ref{sec:operator-algebra} tabulates the full operator algebra.

%--------------------------------------------------------------------
\section{The Composer Operator — Building Larger Resonant Structures}
\label{sec:composer-operator}

The single most powerful construct in the \RCM\ language is the \emph{composer operator} \(\Composer\{\cdot\}\). If an oscillation is a “syllable,” \(\Composer\) is the grammar rule that glues syllables into words and sentences. This section spells out its formal properties, proves its core theorems, and shows how you use it to assemble real systems from small, pitch-matched pieces.

\subsubsection{What \(\Composer\{\cdot\}\) does}

\begin{itemize}[leftmargin=1.5em]
  \item \textbf{Input:} Two (or more) resonance bundles, each defined on a spiral segment.
  \item \textbf{Action:} Attaches the \emph{exit} edge of the first bundle to the \emph{entry} edge of the second, provided the pitches match at the interface.
  \item \textbf{Output:} A new bundle whose spiral path is the geometric concatenation of the inputs, with a composite memory kernel that remembers both histories.
\end{itemize}

\subsubsection{Two governing axioms}

\begin{table}[H]
  \small
  \centering
  \caption{Axioms of the composer operator.}
  \label{tab:composer-axioms}
  \renewcommand{\arraystretch}{1.15}
  \begin{tabularx}{\linewidth}{@{}p{0.25\linewidth} p{0.40\linewidth} X@{}}
    \toprule
    \textbf{Axiom}         & \textbf{Formal statement}                                                        & \textbf{Physical intuition} \\
    \midrule
    Pitch matching         & \(\Composer\{A,B\}\) exists \emph{iff} \(p_{\mathrm{out}}(A)=p_{\mathrm{in}}(B)\) & Coherence passes only when winding rates align. \\
    Associativity          & \(\Composer\{\Composer\{A,B\},C\}=\Composer\{A,\Composer\{B,C\}\}\)             & Gluing order does not affect the final spiral. \\
    \bottomrule
  \end{tabularx}
\end{table}

\noindent\emph{Sketch of proof for associativity:}
Both sides build the same ordered list of spiral segments; since pitch is matched at every internal joint, no re-parameterization is required. Therefore the composite path and its kernel coincide.

\subsubsection{Normal-form theorem}

\textbf{Theorem.}
Any finite nested expression built from \(\Composer\{\cdot\}\) and bundles reduces uniquely to a single bundle whose pitch function is piece-wise continuous and whose memory kernel is a block-convolution of the original kernels.

\paragraph{Idea of proof.}
(1) Use associativity to remove parentheses, obtaining a left-to-right list. (2) Recursively glue adjacent pairs; pitch-continuity (Axiom~1) ensures each step is well-defined. (3) A finite list yields one bundle; uniqueness follows since all parenthesizations reduce to the same list.

\subsubsection{Convolution of memory kernels}

When gluing bundles \(A\) and \(B\), the composite kernel \(K_C\) satisfies
\[
  K_C(\tau)
  = K_A(\tau)
    + \int_{0}^{\tau}
      \bigl(K_A(\tau - \sigma)\searrow K_B(\sigma)\bigr)\,d\sigma.
\]
If \(K_A\) and \(K_B\) are causal exponentials, \(K_C\) becomes a difference of exponentials.

\subsubsection{Worked examples}

\begin{table}[H]
  \small
  \centering
  \caption{Composer operator in practice.}
  \label{tab:composer-examples}
  \renewcommand{\arraystretch}{1.15}
  \begin{tabularx}{\linewidth}{@{}p{0.25\linewidth} p{0.30\linewidth} X@{}}
    \toprule
    \textbf{Task}                  & \textbf{Input bundles}                                               & \textbf{Composer result} \\
    \midrule
    Tight core + loose halo        & \(\RB{\mathrm{core}}\) (pitch \(5^\circ\)), \(\RB{\mathrm{halo}}\) (pitch \(21^\circ\)) &
      Rescale halo by \(T_{\vartheta}\) so its entry pitch becomes \(5^\circ\), then \(\Composer\{\RB{\mathrm{core}},T_{\vartheta}\RB{\mathrm{halo}}\}\). \\
    Theta–gamma phase code         & \(\RB{\theta}\) (8 Hz), \(\RB{\gamma}\) (40 Hz)                       &
      Use kernel enforcing \(5\,\nu_{\theta}=\nu_{\gamma}\); \(\Composer\{\RB{\theta},\RB{\gamma}\}\) carries 5 gamma cycles per theta envelope. \\
    Ring-down echo train           & Initial QNM \(\RB{0}\), back-scatter \(\RB{\mathrm{echo}}\)           &
      \(\Composer\{\RB{0},\RB{\mathrm{echo}},T_{\Delta t}\RB{\mathrm{echo}},T_{2\Delta t}\RB{\mathrm{echo}},\dots\}\). Normal form reduces to one bundle reproducing the full echo sequence. \\
    \bottomrule
  \end{tabularx}
\end{table}

\subsubsection{Category-theory view}

Objects are bundles, morphisms are pitch-matched gluing maps, and \(\Composer\) is the \emph{monoidal product} in the category \textbf{ResBund}. Associativity makes it a strict monoidal structure; pitch-matching is the coherence condition.

\subsubsection{Why \(\Composer\{\cdot\}\) matters in practice}

\begin{itemize}[leftmargin=1.5em]
  \item \textbf{Hierarchical modelling} — Build a galaxy from arms, an arm from clouds, a cloud from scroll waves, using the same operator at each level.
  \item \textbf{Model reuse} — A validated bundle (e.g.\ a neuron) can drop into a larger model (a cortical column) unchanged, provided pitches match.
  \item \textbf{Automated compilation} — A compiler can flatten any nested \(\Composer\) tree to normal form, substitute numeric kernels, and run time-domain simulations with guaranteed continuity at every interface.
\end{itemize}

%--------------------------------------------------------------------
\section{Operator Algebra — How We Transform Resonance Bundles}
\label{sec:operator-algebra}

Section~\ref{sec:composer-operator} fixed the \textbf{glue} (\(\Composer\{\cdot\}\)). Now we catalogue the handful of operators that act on a bundle, prove their algebraic rules, and illustrate how every transformation in later chapters reduces to these primitives.

\begin{table}[H]
  \small
  \centering
  \caption{Core operators in \RCM.}
  \label{tab:rcm-operator-algebra}
  \renewcommand{\arraystretch}{1.15}
  \begin{tabularx}{\linewidth}{@{}p{0.18\linewidth} p{0.20\linewidth} p{0.32\linewidth} X@{}}
    \toprule
    \textbf{Symbol}              & \textbf{Name}               & \textbf{Action on bundle \(\RB{}\)}                     & \textbf{Meaning} \\
    \midrule
    \(\hat{\#}_{\Delta\phi}\)    & Phase shift                 & \(\varsigma(t)\mapsto\varsigma(t)+\Delta\phi\)          & Rotate phase by \(\Delta\phi\). \\
    \(\hat{D}\)                  & Pitch derivative            & Multiplies local amplitude by \(p(\vartheta)\)          & Weight by curvature. \\
    \(\hat{\$}_{n}\)             & Topological-charge projector& Keeps only oscillations with charge \(n\)               & Harmonic/winding-number filter. \\
    \(\hat{M}_{\tau}\)           & Memory filter               & \(A(t)\mapsto\frac1\tau\int_{t-\tau}^{t}A(t')\,dt'\)    & Time-average over past \(\tau\). \\
    \(T_{\vartheta}\)            & Scaling (Lorentz-like)      & \((A,\omega)\mapsto(e^{\vartheta}A,\omega e^{-\vartheta})\) & Rescale space vs.\ time. \\
    \bottomrule
  \end{tabularx}
\end{table}

\subsubsection{Linearity}

All five operators \(O\) are linear on the vector space of bundle amplitudes:
\[
  O\bigl(\alpha\,\RB{1} + \beta\,\RB{2}\bigr)
  = \alpha\,O\,\RB{1} + \beta\,O\,\RB{2}.
\]
\paragraph{Proof sketch.} Each definition is algebraic (shift, multiply) or an integral with a constant kernel; linearity follows immediately.

\subsubsection{Commutators and identities}

\begin{table}[H]
  \small
  \centering
  \caption{Algebraic relations among operators.}
  \label{tab:operator-relations}
  \renewcommand{\arraystretch}{1.15}
  \begin{tabularx}{\linewidth}{@{}p{0.40\linewidth} p{0.30\linewidth} X@{}}
    \toprule
    \textbf{Relation}                                & \textbf{Proof sketch}                                                                                    & \textbf{Interpretation} \\
    \midrule
    \([\hat{\#},\hat{D}]=0\)                         & Phase shift and multiply commute since they act on different “coordinates.”                              & Phase shift doesn’t alter curvature weighting. \\
    \(\hat{\$}_{n}^{2}=\hat{\$}_{n}\)                & Projector applied twice leaves the same harmonic subspace.                                               & Pure harmonic isolation. \\
    \([\hat{D},\hat{M}_{\tau}]\neq 0\)               & Derivative then average \(\neq\) average then derivative unless \(p\) is constant on \([t-\tau,t]\).    & Curvature weighting and averaging do not commute. \\
    \(T_{\vartheta_{1}}T_{\vartheta_{2}}
       = T_{\vartheta_{1}+\vartheta_{2}}\)           & Exponentials add: \(e^{\vartheta_1}e^{\vartheta_2}=e^{\vartheta_1+\vartheta_2}\).                        & \(T\) forms a one-parameter Abelian group. \\
    \bottomrule
  \end{tabularx}
\end{table}

\paragraph{Proof of \([\hat{D},\hat{M}_{\tau}]\neq 0\).}
For any amplitude \(A(t)\):
\[
\begin{aligned}
  \hat{D}\hat{M}_{\tau}A
  &= p(t)\,\frac1\tau\int_{t-\tau}^{t}A(t')\,dt', \\
  \hat{M}_{\tau}\hat{D}A
  &= \frac1\tau\int_{t-\tau}^{t}p(t')A(t')\,dt'.
\end{aligned}
\]
Unless \(p\) is constant on \([t-\tau,t]\), these differ.

\subsubsection{Group and algebra structure}

\begin{itemize}[leftmargin=1.5em]
  \item \(\hat{\#}_{\phi}\) form \(U(1)\): \(\hat{\#}_{\phi_{1}}\hat{\#}_{\phi_{2}}=\hat{\#}_{\phi_{1}+\phi_{2}}\).
  \item \(T_{\vartheta}\) form \((\mathbb{R},+)\), isomorphic to Lorentz rapidity addition.
  \item \(\hat{D}\) acts as the generator of dilations in the emergent Weyl gauge field.
  \item Projectors \(\{\hat{\$}_{n}\}\) are idempotents spanning a Boolean algebra under operator multiplication and addition.
  \item \(\hat{M}_{\tau}\) form a semigroup: \(\hat{M}_{\tau_{1}}\hat{M}_{\tau_{2}}=\hat{M}_{\tau_{1}+\tau_{2}}\).
\end{itemize}

\subsubsection{Normal-ordering prescription}

To resolve ambiguities, adopt the canonical order when multiple operators act:
\[
  \hat{\$}_{n}
  \;\to\;
  \hat{D}
  \;\to\;
  \hat{\#}_{\Delta\phi}
  \;\to\;
  T_{\vartheta}
  \;\to\;
  \hat{M}_{\tau}.
\]
Commutative pairs may swap; non-commuting pairs must follow this sequence for a unique normal form.

\subsubsection{Worked operator chains}

\begin{table}[H]
  \small
  \centering
  \caption{Examples of normal-ordered operator chains.}
  \label{tab:operator-chains}
  \renewcommand{\arraystretch}{1.15}
  \begin{tabularx}{\linewidth}{@{}p{0.35\linewidth} p{0.35\linewidth} X@{}}
    \toprule
    \textbf{Operation}                                        & \textbf{Normal-ordered form}                      & \textbf{Outcome} \\
    \midrule
    Apply \(90^\circ\) phase, then 5 ms filter                & \(\hat{M}_{5\mathrm{ms}}\hat{\#}_{\pi/2}\)        & 5 ms average of a phase-shifted bundle \\
    Boost then pitch derivative                               & \(\hat{\#}_{\Delta\phi}\,\hat{D}\,T_{\vartheta}\) & Curvature-weight, phase-shift, then Lorentz-like boost \\
    Select winding-2 after averaging                          & \(\hat{\$}_{2}\,\hat{M}_{\tau}\)                  & Retain only the 2-winding component of the averaged bundle \\
    \bottomrule
  \end{tabularx}
\end{table}

\subsubsection{Physical examples}

\begin{itemize}[leftmargin=1.5em]
  \item \textbf{Laser cavity alignment} — \(T_{\vartheta}\) models mirror motion; \(\hat{\#}_{\phi}\) adjusts standing-wave nodes; \(\hat{D}\) weights gain profile by curvature for mode-locking.
  \item \textbf{Cortical oscillations} — \(\hat{M}_{100\mathrm{ms}}\) smooths gamma bursts; \(\hat{\$}_{5}\) extracts 5-cycle phase-code within theta.
  \item \textbf{Gravitational-wave post-processing} — boosts \(T_{\vartheta_i}\) account for cosmological redshift; projectors isolate dominant quasinormal modes.
\end{itemize}

\paragraph{Take-aways for §4.3}
\begin{enumerate}[leftmargin=1.5em,label=\arabic*.]
  \item Five core operators generate every local transformation in \RCM.
  \item Their algebra is small: one zero commutator, one nonzero commutator, one idempotence, plus two simple groups.
  \item Normal ordering guarantees uniqueness, enabling unambiguous symbolic manipulation.
  \item Type checking (Term vs.\ Bundle\(\langle k\rangle\)) ensures semantic coherence of compositions.
\end{enumerate}

With the operator algebra catalogued, the next section presents an algorithm that reduces \emph{any} symbolic \RCM\ expression—no matter how nested—to canonical kernel form ready for numerical simulation or analytic proof.

%--------------------------------------------------------------------
\section{Normal-Form Reduction: Canonical Resonance Expressions}

The \RCM\ symbolic framework is powerful but can produce complex, nested expressions. To ensure clarity, consistency, and computational feasibility, every resonance expression is reduced to a canonical \emph{normal form}. This section defines the procedure and theorems that guarantee any valid \RCM\ expression can be systematically simplified.

\subsection{Normal-form reduction — organizing complexity into canonical form}

\begin{quote}
“Complexity without structure is chaos. Normal-form reduction transforms chaotic symbolic complexity into structured, meaningful resonance.”
\end{quote}

\subsubsection{Why normal form?}

Expressions in \RCM\ naturally become nested, especially when modeling realistic physical, biological, or cosmological systems. Without a clear organizing principle, two researchers could write different-looking expressions representing identical physics. Normal-form reduction eliminates ambiguity by defining a canonical form.

\emph{Advantages:}
\begin{itemize}[leftmargin=1.5em]
  \item \textbf{Uniqueness} — Any valid expression reduces to exactly one canonical form.
  \item \textbf{Comparability} — Canonical forms are directly comparable for theoretical or computational checks.
  \item \textbf{Automation} — Enables symbolic tools to simplify complex resonant expressions reliably.
\end{itemize}

\subsubsection{The canonical operator order}

We fix the operator order
\[
  \hat{\$}_n \;\to\; \hat{D} \;\to\; \hat{\#}_{\phi} \;\to\; T_{\vartheta} \;\to\; \hat{M}_{\tau},
\]
for the reasons given in Section~\ref{sec:operator-algebra}.

\subsubsection{Rewrite rules: achieving normal form}

The following terminating and confluent rules transform any expression to canonical form:

\begin{table}[H]
  \small
  \centering
  \caption{Rewrite rules for normal-form reduction.}
  \label{tab:rewrite-rules}
  \renewcommand{\arraystretch}{1.15}
  \begin{tabularx}{\linewidth}{@{}p{0.30\linewidth} p{0.30\linewidth} X@{}}
    \toprule
    \textbf{Original pattern}                  & \textbf{Rewrite}                          & \textbf{Explanation} \\
    \midrule
    \(\Composer\{\Composer\{A,B\},C\}\)        & \(\Composer\{A,B,C\}\)                    & Flatten nested composers (associativity). \\
    \(T_{\vartheta_1}T_{\vartheta_2}\)         & \(T_{\vartheta_1+\vartheta_2}\)           & Combine successive scalings. \\
    \(\hat{\#}_{\phi_1}\hat{\#}_{\phi_2}\)     & \(\hat{\#}_{\phi_1+\phi_2}\)              & Merge phase shifts. \\
    \(\hat{\$}_n\hat{\$}_n\)                   & \(\hat{\$}_n\)                            & Idempotence of projectors. \\
    \(\hat{\$}_n\hat{\$}_m,\,n\neq m\)         & \(0\)                                     & Orthogonal projections. \\
    \(\hat{D}\hat{\#}_{\phi}\)                 & \(\hat{\#}_{\phi}\hat{D}\)                & Commutation. \\
    \(\hat{D}\hat{M}_{\tau}\)                  & \(\hat{M}_{\tau}\hat{D}+[\hat{D},\hat{M}_{\tau}]\) & Non-commutation made explicit. \\
    \(\hat{M}_{\tau_1}\hat{M}_{\tau_2}\)       & \(\hat{M}_{\tau_1+\tau_2}\)               & Filters compose. \\
    \bottomrule
  \end{tabularx}
\end{table}

\subsubsection{Algorithmic procedure (step-by-step)}

For an arbitrary \RCM\ expression \(E\):

\begin{enumerate}[leftmargin=1.5em]
  \item \textbf{Flatten composers:} \(\Composer\{\Composer\{A,B\},C\}\to\Composer\{A,B,C\}\).
  \item \textbf{Match boundary pitches:} insert minimal \(T_{\vartheta}\) to enforce \(p_{\mathrm{out}}(A)=p_{\mathrm{in}}(B)\).
  \item \textbf{Combine similar operators:} apply Table~\ref{tab:rewrite-rules}.
  \item \textbf{Reorder to canonical form:} follow the fixed operator order; introduce commutators when needed.
  \item \textbf{Collapse kernels:} \(K_{\mathrm{total}}=K_{1}\searrow K_{2}\searrow\cdots\searrow K_{n}\).
\end{enumerate}

Repeat until no further simplifications apply.

\subsubsection{Concrete example of normal-form reduction}

Consider
\[
  E = \hat{M}_{10\mathrm{ms}}\,
      \hat{\#}_{\pi/6}\,
      \hat{D}\,
      T_{0.2}\,
      \hat{\$}_{2}\,
      \Composer\bigl\{\Composer\{\RB{a},\RB{b}\},\,T_{-0.2}\hat{\#}_{-\pi/6}\RB{c}\bigr\}.
\]
\begin{enumerate}[leftmargin=1.5em]
  \item Flatten: \(\Composer\{\RB{a},\RB{b},T_{-0.2}\hat{\#}_{-\pi/6}\RB{c}\}\).
  \item Pitch-match: insert \(T_{0.2}\) at the interface.
  \item Collapse scaling: \(T_{0.2}T_{-0.2}\to I\).
  \item Combine phase shifts: \(\hat{\#}_{\pi/6}\hat{\#}_{-\pi/6}\to I\).
  \item Canonical order: \(\hat{\$}_{2}\to \hat{D}\to \hat{M}_{10\mathrm{ms}}\).
\end{enumerate}
Thus
\[
  E_{\mathrm{norm}}
  = \hat{\$}_{2}\,\hat{D}\,\hat{M}_{10\mathrm{ms}}\,\RB{ABC},
  \qquad
  \RB{ABC}=\Composer\{\RB{a},\RB{b},\RB{c}\},
  \quad
  K_{ABC} = K_a\searrow K_b\searrow K_c.
\]

\subsubsection{Summary and link to aggregation matrix}

Normal-form reduction ensures every \RCM\ expression is uniquely defined, computationally practical, and conceptually transparent. With expressions canonicalized, we assemble these kernels into the global \emph{aggregation matrix}, encoding the collective resonance behavior of entire systems.

%--------------------------------------------------------------------
\section{Kernel Functions and Aggregation Matrix}

\begin{center}
\emph{“Memory is the resonance of past interactions; the kernel is the mathematical trace they leave behind.”}
\end{center}

\subsubsection{Memory kernel: the mathematical signature of interaction}

Every encounter between resonance bundles leaves a mathematical signature: the \textbf{memory kernel} \(K_{\alpha\beta}(\tau)\), representing how bundle \(\Psi_\beta\) influences bundle \(\Psi_\alpha\) with delay \(\tau\):
\[
  K_{\alpha\beta}(\tau)=0 \quad \text{for } \tau<0.
\]
The bundle’s evolution is then influenced explicitly by its past interactions:
\[
  \dot{\Psi}_\alpha(t)
  = F_\alpha[\Psi_\alpha(t)]
  + \sum_{\beta}\int_{-\infty}^{t} K_{\alpha\beta}(t-t')\,\Psi_\beta(t')\,dt'.
\]

\subsubsection{Aggregation matrix: weaving memory into resonance}

Define the \textbf{aggregation matrix}
\[
  M_{\alpha\beta}
  = \iint K_{\alpha\beta}(t-t')\,\Psi_\alpha^*(t)\,\Psi_\beta(t')\,dt\,dt'.
\]
It encodes which resonant states are supported by the system’s collective memory.

\subsubsection{Resonance stability as an emergent property}

Stability arises from the spectral properties of \(M\):
\[
  M\,\mathbf{v}_n = \lambda_n \mathbf{v}_n,
\]
with \(\Re(\lambda_n)\) determining growth/decay and \(\Im(\lambda_n)\) setting oscillation rates.

\subsubsection{Symmetry, conservation, and implications}

Symmetries in kernels imply conservation laws, e.g.
\[
  \frac{d}{dt}\sum_{\alpha}\langle\Psi_\alpha|\Psi_\alpha\rangle=0,
\]
reflecting invariances across scales.

\subsubsection{Numerical computation: from philosophy to practice}

\begin{enumerate}[leftmargin=1.5em]
  \item Discretize memory integrals:
  \[
    M_{\alpha\beta}\approx\sum_{ij}K_{\alpha\beta}(t_i-t_j)\,\Psi_\alpha^*(t_i)\,\Psi_\beta(t_j)\,\Delta t^2.
  \]
  \item Extract resonant structures via eigendecomposition:
  \[
    M=V\Lambda V^{-1},\qquad \Lambda=\mathrm{diag}(\lambda_n).
  \]
\end{enumerate}

\subsubsection{Physical manifestations: memory in nature}

\begin{itemize}[leftmargin=1.5em]
  \item \textbf{Optical resonance}: cavity modes “remember” boundary reflections.
  \item \textbf{Biological rhythms}: neural oscillations encode memory of past firings.
  \item \textbf{Cosmological echoes}: gravitational waves remember black-hole mergers as late-time tails.
\end{itemize}
In each context, resonance is memory made manifest.
